\documentclass[11pt]{article}
\newcommand{\ListaTitulo}{Gabarito 11: VA contínua, esperança, variância e FDA}
% Preâmbulo compartilhado: MATF14, Listas de Exercícios 2026
% Usado por ListaNN.tex e GabaritoNN.tex via % Preâmbulo compartilhado: MATF14, Listas de Exercícios 2026
% Usado por ListaNN.tex e GabaritoNN.tex via % Preâmbulo compartilhado: MATF14, Listas de Exercícios 2026
% Usado por ListaNN.tex e GabaritoNN.tex via \input{preamble}

\usepackage[utf8]{inputenc}
\usepackage[T1]{fontenc}
\usepackage[brazil]{babel}
\usepackage{fullpage}
\usepackage{amsmath,amssymb}
\usepackage{enumitem}
\usepackage{xcolor}
\usepackage{fancyhdr}
\usepackage{titlesec}
\usepackage{listings}
\usepackage{hyperref}
\usepackage{booktabs}
\usepackage{graphicx}
\usepackage{tikz}

\definecolor{matf14azul}{HTML}{0D3B54}
\definecolor{matf14dourado}{HTML}{C9971E}

\titleformat{\section}{\normalfont\Large\bfseries\color{matf14azul}}{\thesection}{1em}{}
\titleformat{\subsection}{\normalfont\large\bfseries\color{matf14azul}}{\thesubsection}{1em}{}

\providecommand{\ListaTitulo}{Lista de Exercícios}

\setlength{\headheight}{14pt}
\pagestyle{fancy}
\fancyhf{}
\lhead{\small MATF14: Estatística Econômica I}
\rhead{\small \ListaTitulo}
\cfoot{\thepage}
\renewcommand{\headrulewidth}{0.4pt}

\lstset{
  basicstyle=\ttfamily\small,
  breaklines=true,
  frame=single,
  columns=fullflexible,
  backgroundcolor=\color{gray!8},
  commentstyle=\color{matf14dourado},
  keywordstyle=\color{matf14azul}\bfseries,
  language=R
}

\newcounter{questaocounter}
\newenvironment{questao}[1]{%
  \stepcounter{questaocounter}%
  \bigskip\noindent\textbf{Questão #1.}\ %
}{\par}

\newcommand{\cabecalho}[2]{%
  \begin{center}
    {\Large\bfseries\color{matf14azul} #1}\\[0.3em]
    {\large #2}
  \end{center}
  \vspace{0.5em}
  \noindent\rule{\linewidth}{0.6pt}
  \vspace{0.5em}
}

\newenvironment{solucao}{%
  \par\smallskip\noindent\textit{\color{matf14dourado!80!black}Solução.}\ %
}{\hfill$\blacksquare$\par}


\usepackage[utf8]{inputenc}
\usepackage[T1]{fontenc}
\usepackage[brazil]{babel}
\usepackage{fullpage}
\usepackage{amsmath,amssymb}
\usepackage{enumitem}
\usepackage{xcolor}
\usepackage{fancyhdr}
\usepackage{titlesec}
\usepackage{listings}
\usepackage{hyperref}
\usepackage{booktabs}
\usepackage{graphicx}
\usepackage{tikz}

\definecolor{matf14azul}{HTML}{0D3B54}
\definecolor{matf14dourado}{HTML}{C9971E}

\titleformat{\section}{\normalfont\Large\bfseries\color{matf14azul}}{\thesection}{1em}{}
\titleformat{\subsection}{\normalfont\large\bfseries\color{matf14azul}}{\thesubsection}{1em}{}

\providecommand{\ListaTitulo}{Lista de Exercícios}

\setlength{\headheight}{14pt}
\pagestyle{fancy}
\fancyhf{}
\lhead{\small MATF14: Estatística Econômica I}
\rhead{\small \ListaTitulo}
\cfoot{\thepage}
\renewcommand{\headrulewidth}{0.4pt}

\lstset{
  basicstyle=\ttfamily\small,
  breaklines=true,
  frame=single,
  columns=fullflexible,
  backgroundcolor=\color{gray!8},
  commentstyle=\color{matf14dourado},
  keywordstyle=\color{matf14azul}\bfseries,
  language=R
}

\newcounter{questaocounter}
\newenvironment{questao}[1]{%
  \stepcounter{questaocounter}%
  \bigskip\noindent\textbf{Questão #1.}\ %
}{\par}

\newcommand{\cabecalho}[2]{%
  \begin{center}
    {\Large\bfseries\color{matf14azul} #1}\\[0.3em]
    {\large #2}
  \end{center}
  \vspace{0.5em}
  \noindent\rule{\linewidth}{0.6pt}
  \vspace{0.5em}
}

\newenvironment{solucao}{%
  \par\smallskip\noindent\textit{\color{matf14dourado!80!black}Solução.}\ %
}{\hfill$\blacksquare$\par}


\usepackage[utf8]{inputenc}
\usepackage[T1]{fontenc}
\usepackage[brazil]{babel}
\usepackage{fullpage}
\usepackage{amsmath,amssymb}
\usepackage{enumitem}
\usepackage{xcolor}
\usepackage{fancyhdr}
\usepackage{titlesec}
\usepackage{listings}
\usepackage{hyperref}
\usepackage{booktabs}
\usepackage{graphicx}
\usepackage{tikz}

\definecolor{matf14azul}{HTML}{0D3B54}
\definecolor{matf14dourado}{HTML}{C9971E}

\titleformat{\section}{\normalfont\Large\bfseries\color{matf14azul}}{\thesection}{1em}{}
\titleformat{\subsection}{\normalfont\large\bfseries\color{matf14azul}}{\thesubsection}{1em}{}

\providecommand{\ListaTitulo}{Lista de Exercícios}

\setlength{\headheight}{14pt}
\pagestyle{fancy}
\fancyhf{}
\lhead{\small MATF14: Estatística Econômica I}
\rhead{\small \ListaTitulo}
\cfoot{\thepage}
\renewcommand{\headrulewidth}{0.4pt}

\lstset{
  basicstyle=\ttfamily\small,
  breaklines=true,
  frame=single,
  columns=fullflexible,
  backgroundcolor=\color{gray!8},
  commentstyle=\color{matf14dourado},
  keywordstyle=\color{matf14azul}\bfseries,
  language=R
}

\newcounter{questaocounter}
\newenvironment{questao}[1]{%
  \stepcounter{questaocounter}%
  \bigskip\noindent\textbf{Questão #1.}\ %
}{\par}

\newcommand{\cabecalho}[2]{%
  \begin{center}
    {\Large\bfseries\color{matf14azul} #1}\\[0.3em]
    {\large #2}
  \end{center}
  \vspace{0.5em}
  \noindent\rule{\linewidth}{0.6pt}
  \vspace{0.5em}
}

\newenvironment{solucao}{%
  \par\smallskip\noindent\textit{\color{matf14dourado!80!black}Solução.}\ %
}{\hfill$\blacksquare$\par}


\begin{document}

\cabecalho{Gabarito 11}{Variável aleatória contínua, densidade, esperança, variância e FDA (Cap. 4, Aulas 24--26)}

\begin{questao}{1} \textbf{(Densidade válida: tempo de vida de um equipamento)}
\begin{solucao}
\begin{enumerate}[label=(\alph*)]
  \item $\int_0^4 kx\,dx = k\left[\frac{x^2}{2}\right]_0^4 = k\times 8 = 1 \Rightarrow k=1/8$.
  \item $F(x) = \int_0^x \frac{t}{8}dt = \frac{x^2}{16}$. $P(X\le2)=F(2)=4/16=0{,}25$.
  \item $P(1\le X\le3) = F(3)-F(1) = 9/16-1/16 = 8/16=0{,}5$.
\end{enumerate}
\end{solucao}
\end{questao}

\begin{questao}{2} \textbf{(Esperança e variância: receita de um posto de vendas)}
\begin{solucao}
\begin{enumerate}[label=(\alph*)]
  \item $E(X)=\int_2^{10} x\cdot\frac18\,dx = \frac18\left[\frac{x^2}{2}\right]_2^{10} =
    \frac18\left(\frac{100-4}{2}\right) = \frac18\times48=6$.
  \item $\mathrm{Var}(X) = \int_2^{10}(x-6)^2\cdot\frac18\,dx$. Fazendo $u=x-6$, limites $-4$ a
    $4$: $\frac18\int_{-4}^4 u^2\,du = \frac18\left[\frac{u^3}{3}\right]_{-4}^4 =
    \frac18\times\frac{64-(-64)}{3} = \frac18\times\frac{128}{3} = \frac{16}{3}\approx5{,}33$.
  \item Fórmula pronta: $E(X)=\frac{2+10}{2}=6$ (\checkmark). $\mathrm{Var}(X)=\frac{(10-2)^2}{12}=
    \frac{64}{12}=\frac{16}{3}\approx5{,}33$ (\checkmark). Ambos coincidem com (a) e (b).
\end{enumerate}
\end{solucao}
\end{questao}

\begin{questao}{3} \textbf{(FDA a partir da densidade: desconto em uma promoção)}
\begin{solucao}
\begin{enumerate}[label=(\alph*)]
  \item $F(x) = \int_0^x \frac{3}{8000}t^2\,dt = \frac{3}{8000}\left[\frac{t^3}{3}\right]_0^x =
    \frac{x^3}{8000}$, para $0\le x\le20$.
  \item $P(X\ge15) = 1-F(15) = 1-\frac{15^3}{8000} = 1-\frac{3375}{8000} = 1-0{,}4219 =
    0{,}5781$.
  \item $P(X\ge15) = \int_{15}^{20}\frac{3}{8000}x^2\,dx = \left[\frac{x^3}{8000}\right]_{15}^{20}
    = \frac{8000}{8000}-\frac{3375}{8000} = 1-0{,}4219=0{,}5781$. Mesmo resultado, como esperado
    (o item (b) é só o Teorema Fundamental do Cálculo aplicado à mesma integral).
\end{enumerate}
\end{solucao}
\end{questao}

\begin{questao}{4} \textbf{(Propriedades: conversão de moeda)}
\begin{solucao}
\begin{enumerate}[label=(\alph*)]
  \item $E(C) = 5{,}20\times50+30 = 260+30=290$.
  \item $dp(C) = |5{,}20|\times dp(X) = 5{,}20\times8=41{,}6$.
  \item O desvio padrão \textbf{aumenta}: $dp(C)=5{,}50\times8=44{,}0 > 41{,}6$. Pela propriedade
    $\mathrm{Var}(aX+b)=a^2\mathrm{Var}(X)$ (logo $dp(aX+b)=|a|\,dp(X)$), o desvio padrão é
    diretamente proporcional ao coeficiente de escala $a$ (a cotação); a taxa fixa $b=30$ não
    entra nessa conta.
\end{enumerate}
\end{solucao}
\end{questao}

\begin{questao}{5} \textbf{(Dedução: $P(X=x)=0$ para VA contínua)}
\begin{solucao}
\begin{enumerate}[label=(\alph*)]
  \item $P(x-\varepsilon<X\le x) = F(x) - F(x-\varepsilon)$ (diferença de duas probabilidades
    acumuladas).
  \item Por definição, $F$ é contínua em $x$ significa exatamente que
    $\lim_{\varepsilon\to0^+} F(x-\varepsilon) = F(x)$, logo
    $\lim_{\varepsilon\to0^+}[F(x)-F(x-\varepsilon)] = F(x)-F(x) = 0$.
  \item O evento $\{X=x\}$ é o limite dos eventos $\{x-\varepsilon<X\le x\}$ quando
    $\varepsilon\to0^+$ (o intervalo encolhe até o ponto único $x$). Pela continuidade de
    probabilidade (um resultado geral que estende os axiomas), $P(X=x) =
    \lim_{\varepsilon\to0^+} P(x-\varepsilon<X\le x)$, que pelo item (b) vale $0$. Se $F$ tivesse
    um salto em $x$ (caso discreto), o limite em (b) seria igual ao tamanho do salto, não zero,
    é exatamente essa ausência de saltos, garantida pela continuidade de $F$, que força
    $P(X=x)=0$ no caso contínuo.
\end{enumerate}
\end{solucao}
\end{questao}

\end{document}
